% --- Archon physics formalization source begin ---
% archon:physics
% archon:covers PhyXMiniProblems/problem_phyx_mini_0879.lean
% archon:source-report reports/phyx_mini/problem_phyx_mini_0879.source.json

\chapter{Physics problem phyx\_mini\_0879}
\label{ch:PhyXMiniProblems_problem_phyx_mini_0879}

\paragraph{Problem source.}
\#\# Question

What is the equal charge on the three capacitors?

\#\# Answer choices

- A: A: 2.5\textbackslash{} \textbackslash{}mu\textbackslash{}mathrm\{C\}

- B: B: 35\textbackslash{} \textbackslash{}mu\textbackslash{}mathrm\{C\}

- C: C: 55\textbackslash{} \textbackslash{}mu\textbackslash{}mathrm\{C\}

- D: D: 60\textbackslash{} \textbackslash{}mu\textbackslash{}mathrm\{C\}

\#\# Figure caption (auxiliary; use the image as primary evidence)

The image is a circuit diagram featuring a voltage source and three capacitors connected in series.

1. **Voltage Source**: On the left side, there is a labeled voltage source of \textbackslash{}(30 \textbackslash{}, \textbackslash{}text\{V\}\textbackslash{}) indicated by a rectangular symbol with a plus and minus sign.

2. **Capacitors**: On the right side of the circuit, there are three capacitors in series. 
   - The first capacitor is labeled \textbackslash{}(C\_1 = 12 \textbackslash{}, \textbackslash{}mu\textbackslash{}text\{F\}\textbackslash{}).
   - The second capacitor is labeled \textbackslash{}(C\_2 = 4 \textbackslash{}, \textbackslash{}mu\textbackslash{}text\{F\}\textbackslash{}).
   - The third capacitor is labeled \textbackslash{}(C\_3 = 6 \textbackslash{}, \textbackslash{}mu\textbackslash{}text\{F\}\textbackslash{}).

3. **Circuit Connections**: The capacitors are connected in a single straight line, indicating they are in series. The voltage source is connected across this series configuration, completing the circuit.

The image shows the relationship between the voltage source and the three capacitors within the circuit.

\#\# Dataset metadata

Electric Circuits; Physical Model Grounding Reasoning, Multi-Formula Reasoning

\paragraph{Recorded answer/context.}
D: D: 60\textbackslash{} \textbackslash{}mu\textbackslash{}mathrm\{C\}

\paragraph{Figure/image path.}
/root/proposal\_for\_physic/science-mango/phyx\_mini\_run/phyx\_data/test\_image/879.png

\paragraph{Formalization target.}
create a compiling Lean file with sorry bodies at `PhyXMiniProblems/problem\_phyx\_mini\_0879.lean`. The Lean declarations must preserve the physical quantities, dimensions or dimensional roles, figure labels, governing-law hypotheses, and final relation expressed by this problem.
Use Mathlib/Physlib names found through LeanExplore where available. If a physics API is missing, introduce faithful local abstractions rather than scalar placeholder aliases.

\begin{theorem}[Physics formalization target]
\label{thm:physics:phyx_mini_0879:target}
\lean{PhyXMiniProblems.ProblemPhyXMini0879.problem_phyx_mini_0879}
\uses{def:physics:phyx-mini-0879:phyxminiproblems-problemphyxmini0879-chargemagnitudeinmicrocoulombs, def:physics:phyx-mini-0879:phyxminiproblems-problemphyxmini0879-threeseriescapacitorsetup, def:physics:phyx-mini-0879:phyxminiproblems-problemphyxmini0879-matchesthreecapacitorseriestopology, def:physics:phyx-mini-0879:phyxminiproblems-problemphyxmini0879-matchessuppliedthreecapacitorfigure, def:physics:phyx-mini-0879:phyxminiproblems-problemphyxmini0879-hasphysicalseriescapacitorparameters, def:physics:phyx-mini-0879:phyxminiproblems-problemphyxmini0879-satisfiesidealcapacitorlaws, def:physics:phyx-mini-0879:phyxminiproblems-problemphyxmini0879-satisfiesseriescircuitconservationlaws, def:physics:phyx-mini-0879:phyxminiproblems-problemphyxmini0879-recordeddatasetanswer, def:physics:phyx-mini-0879:phyxminiproblems-problemphyxmini0879-isuniquematchinganswer}
The assigned autoformalize agent should translate this physics problem into Lean declarations in the covered file, with theorem and lemma proof bodies written as `by sorry`.
\end{theorem}
\begin{proof}
This is an autoformalization task, not a proof task. Produce statements that can later be proved by the physics prover without weakening the physical model.
\end{proof}

% --- Archon declaration topology begin ---
\section{Declaration topology}
\begin{definition}[potential Difference Dimension]
  \label{def:physics:phyx-mini-0879:phyxminiproblems-problemphyxmini0879-potentialdifferencedimension}
  \lean{PhyXMiniProblems.ProblemPhyXMini0879.potentialDifferenceDimension}
  The dimension M L² T⁻² C⁻¹ of electric potential difference
\end{definition}
\begin{proof}
  This is the declared model definition of potential Difference Dimension.
\end{proof}
\begin{definition}[capacitance Dimension]
  \label{def:physics:phyx-mini-0879:phyxminiproblems-problemphyxmini0879-capacitancedimension}
  \lean{PhyXMiniProblems.ProblemPhyXMini0879.capacitanceDimension}
  \uses{def:physics:phyx-mini-0879:phyxminiproblems-problemphyxmini0879-potentialdifferencedimension}
  The dimension C / V of electrical capacitance
\end{definition}
\begin{proof}
  This is the declared model definition of capacitance Dimension.
\end{proof}
\begin{definition}[Potential Difference Quantity]
  \label{def:physics:phyx-mini-0879:phyxminiproblems-problemphyxmini0879-potentialdifferencequantity}
  \lean{PhyXMiniProblems.ProblemPhyXMini0879.PotentialDifferenceQuantity}
  \uses{def:physics:phyx-mini-0879:phyxminiproblems-problemphyxmini0879-potentialdifferencedimension}
  A nonnegative, unit-independent potential-difference magnitude
\end{definition}
\begin{proof}
  This is the declared model definition of Potential Difference Quantity.
\end{proof}
\begin{definition}[Capacitance Quantity]
  \label{def:physics:phyx-mini-0879:phyxminiproblems-problemphyxmini0879-capacitancequantity}
  \lean{PhyXMiniProblems.ProblemPhyXMini0879.CapacitanceQuantity}
  \uses{def:physics:phyx-mini-0879:phyxminiproblems-problemphyxmini0879-capacitancedimension}
  A nonnegative, unit-independent electrical capacitance
\end{definition}
\begin{proof}
  This is the declared model definition of Capacitance Quantity.
\end{proof}
\begin{definition}[Charge Magnitude Quantity]
  \label{def:physics:phyx-mini-0879:phyxminiproblems-problemphyxmini0879-chargemagnitudequantity}
  \lean{PhyXMiniProblems.ProblemPhyXMini0879.ChargeMagnitudeQuantity}
  A nonnegative, unit-independent electric-charge magnitude
\end{definition}
\begin{proof}
  This is the declared model definition of Charge Magnitude Quantity.
\end{proof}
\begin{definition}[Signed Charge Quantity]
  \label{def:physics:phyx-mini-0879:phyxminiproblems-problemphyxmini0879-signedchargequantity}
  \lean{PhyXMiniProblems.ProblemPhyXMini0879.SignedChargeQuantity}
  A signed, unit-independent electric charge on one capacitor plate
\end{definition}
\begin{proof}
  This is the declared model definition of Signed Charge Quantity.
\end{proof}
\begin{definition}[potential Difference In Volts]
  \label{def:physics:phyx-mini-0879:phyxminiproblems-problemphyxmini0879-potentialdifferenceinvolts}
  \lean{PhyXMiniProblems.ProblemPhyXMini0879.potentialDifferenceInVolts}
  \uses{def:physics:phyx-mini-0879:phyxminiproblems-problemphyxmini0879-potentialdifferencequantity}
  Read a potential-difference magnitude in coherent-SI volts
\end{definition}
\begin{proof}
  This is the declared model definition of potential Difference In Volts.
\end{proof}
\begin{definition}[capacitance In Farads]
  \label{def:physics:phyx-mini-0879:phyxminiproblems-problemphyxmini0879-capacitanceinfarads}
  \lean{PhyXMiniProblems.ProblemPhyXMini0879.capacitanceInFarads}
  \uses{def:physics:phyx-mini-0879:phyxminiproblems-problemphyxmini0879-capacitancequantity}
  Read an electrical capacitance in coherent-SI farads
\end{definition}
\begin{proof}
  This is the declared model definition of capacitance In Farads.
\end{proof}
\begin{definition}[capacitance In Microfarads]
  \label{def:physics:phyx-mini-0879:phyxminiproblems-problemphyxmini0879-capacitanceinmicrofarads}
  \lean{PhyXMiniProblems.ProblemPhyXMini0879.capacitanceInMicrofarads}
  \uses{def:physics:phyx-mini-0879:phyxminiproblems-problemphyxmini0879-capacitancequantity, def:physics:phyx-mini-0879:phyxminiproblems-problemphyxmini0879-capacitanceinfarads}
  Read an electrical capacitance in microfarads
\end{definition}
\begin{proof}
  This is the declared model definition of capacitance In Microfarads.
\end{proof}
\begin{definition}[micro Coulombs]
  \label{def:physics:phyx-mini-0879:phyxminiproblems-problemphyxmini0879-microcoulombs}
  \lean{PhyXMiniProblems.ProblemPhyXMini0879.microCoulombs}
  The charge unit 10⁻⁶ C used by the requested answer
\end{definition}
\begin{proof}
  This is the declared model definition of micro Coulombs.
\end{proof}
\begin{definition}[charge Magnitude Readout]
  \label{def:physics:phyx-mini-0879:phyxminiproblems-problemphyxmini0879-chargemagnitudereadout}
  \lean{PhyXMiniProblems.ProblemPhyXMini0879.chargeMagnitudeReadout}
  \uses{def:physics:phyx-mini-0879:phyxminiproblems-problemphyxmini0879-chargemagnitudequantity}
  Read a charge magnitude in a selected physical charge unit
\end{definition}
\begin{proof}
  This is the declared model definition of charge Magnitude Readout.
\end{proof}
\begin{definition}[signed Charge Readout]
  \label{def:physics:phyx-mini-0879:phyxminiproblems-problemphyxmini0879-signedchargereadout}
  \lean{PhyXMiniProblems.ProblemPhyXMini0879.signedChargeReadout}
  \uses{def:physics:phyx-mini-0879:phyxminiproblems-problemphyxmini0879-signedchargequantity}
  Read a signed plate charge in a selected physical charge unit
\end{definition}
\begin{proof}
  This is the declared model definition of signed Charge Readout.
\end{proof}
\begin{definition}[charge Magnitude In Coulombs]
  \label{def:physics:phyx-mini-0879:phyxminiproblems-problemphyxmini0879-chargemagnitudeincoulombs}
  \lean{PhyXMiniProblems.ProblemPhyXMini0879.chargeMagnitudeInCoulombs}
  \uses{def:physics:phyx-mini-0879:phyxminiproblems-problemphyxmini0879-chargemagnitudequantity, def:physics:phyx-mini-0879:phyxminiproblems-problemphyxmini0879-chargemagnitudereadout}
  Read a capacitor's charge magnitude in coulombs
\end{definition}
\begin{proof}
  This is the declared model definition of charge Magnitude In Coulombs.
\end{proof}
\begin{definition}[charge Magnitude In Microcoulombs]
  \label{def:physics:phyx-mini-0879:phyxminiproblems-problemphyxmini0879-chargemagnitudeinmicrocoulombs}
  \lean{PhyXMiniProblems.ProblemPhyXMini0879.chargeMagnitudeInMicrocoulombs}
  \uses{def:physics:phyx-mini-0879:phyxminiproblems-problemphyxmini0879-chargemagnitudequantity, def:physics:phyx-mini-0879:phyxminiproblems-problemphyxmini0879-microcoulombs, def:physics:phyx-mini-0879:phyxminiproblems-problemphyxmini0879-chargemagnitudereadout}
  Read a capacitor's charge magnitude in microcoulombs
\end{definition}
\begin{proof}
  This is the declared model definition of charge Magnitude In Microcoulombs.
\end{proof}
\begin{definition}[signed Charge In Coulombs]
  \label{def:physics:phyx-mini-0879:phyxminiproblems-problemphyxmini0879-signedchargeincoulombs}
  \lean{PhyXMiniProblems.ProblemPhyXMini0879.signedChargeInCoulombs}
  \uses{def:physics:phyx-mini-0879:phyxminiproblems-problemphyxmini0879-signedchargequantity, def:physics:phyx-mini-0879:phyxminiproblems-problemphyxmini0879-signedchargereadout}
  Read a signed capacitor-plate charge in coulombs
\end{definition}
\begin{proof}
  This is the declared model definition of signed Charge In Coulombs.
\end{proof}
\begin{definition}[Capacitor Label]
  \label{def:physics:phyx-mini-0879:phyxminiproblems-problemphyxmini0879-capacitorlabel}
  \lean{PhyXMiniProblems.ProblemPhyXMini0879.CapacitorLabel}
  The three capacitors, named as in the supplied image
\end{definition}
\begin{proof}
  This is the declared model definition of Capacitor Label.
\end{proof}
\begin{definition}[Capacitor Plate]
  \label{def:physics:phyx-mini-0879:phyxminiproblems-problemphyxmini0879-capacitorplate}
  \lean{PhyXMiniProblems.ProblemPhyXMini0879.CapacitorPlate}
  The two plates of a capacitor in the image's top-to-bottom orientation
\end{definition}
\begin{proof}
  This is the declared model definition of Capacitor Plate.
\end{proof}
\begin{definition}[Circuit Node]
  \label{def:physics:phyx-mini-0879:phyxminiproblems-problemphyxmini0879-circuitnode}
  \lean{PhyXMiniProblems.ProblemPhyXMini0879.CircuitNode}
  Electrical nodes in the single closed series loop
\end{definition}
\begin{proof}
  This is the declared model definition of Circuit Node.
\end{proof}
\begin{definition}[Capacitor Circuit Model]
  \label{def:physics:phyx-mini-0879:phyxminiproblems-problemphyxmini0879-capacitorcircuitmodel}
  \lean{PhyXMiniProblems.ProblemPhyXMini0879.CapacitorCircuitModel}
  The ideal circuit interpretation of the diagram
\end{definition}
\begin{proof}
  This is the declared model definition of Capacitor Circuit Model.
\end{proof}
\begin{definition}[Three Series Capacitor Figure]
  \label{def:physics:phyx-mini-0879:phyxminiproblems-problemphyxmini0879-threeseriescapacitorfigure}
  \lean{PhyXMiniProblems.ProblemPhyXMini0879.ThreeSeriesCapacitorFigure}
  \uses{def:physics:phyx-mini-0879:phyxminiproblems-problemphyxmini0879-capacitorlabel}
  Literal visual information carried by image 879.png
\end{definition}
\begin{proof}
  This is the declared model definition of Three Series Capacitor Figure.
\end{proof}
\begin{definition}[Three Series Capacitor Setup]
  \label{def:physics:phyx-mini-0879:phyxminiproblems-problemphyxmini0879-threeseriescapacitorsetup}
  \lean{PhyXMiniProblems.ProblemPhyXMini0879.ThreeSeriesCapacitorSetup}
  \uses{def:physics:phyx-mini-0879:phyxminiproblems-problemphyxmini0879-potentialdifferencequantity, def:physics:phyx-mini-0879:phyxminiproblems-problemphyxmini0879-capacitancequantity, def:physics:phyx-mini-0879:phyxminiproblems-problemphyxmini0879-chargemagnitudequantity, def:physics:phyx-mini-0879:phyxminiproblems-problemphyxmini0879-signedchargequantity, def:physics:phyx-mini-0879:phyxminiproblems-problemphyxmini0879-capacitorlabel, def:physics:phyx-mini-0879:phyxminiproblems-problemphyxmini0879-capacitorplate, def:physics:phyx-mini-0879:phyxminiproblems-problemphyxmini0879-circuitnode, def:physics:phyx-mini-0879:phyxminiproblems-problemphyxmini0879-capacitorcircuitmodel, def:physics:phyx-mini-0879:phyxminiproblems-problemphyxmini0879-threeseriescapacitorfigure}
  The independent physical data of the circuit. The node maps retain the figure's connection topology. Charge magnitudes and signed plate charges are not defined from the answer or from one another; the general capacitor and charge-conservation laws relate them below
\end{definition}
\begin{proof}
  This is the declared model definition of Three Series Capacitor Setup.
\end{proof}
\begin{definition}[Matches Three Capacitor Series Topology]
  \label{def:physics:phyx-mini-0879:phyxminiproblems-problemphyxmini0879-matchesthreecapacitorseriestopology}
  \lean{PhyXMiniProblems.ProblemPhyXMini0879.MatchesThreeCapacitorSeriesTopology}
  \uses{def:physics:phyx-mini-0879:phyxminiproblems-problemphyxmini0879-threeseriescapacitorsetup}
  The single-branch series topology represented by the pictured wires
\end{definition}
\begin{proof}
  This is the declared model definition of Matches Three Capacitor Series Topology.
\end{proof}
\begin{definition}[Matches Supplied Three Capacitor Figure]
  \label{def:physics:phyx-mini-0879:phyxminiproblems-problemphyxmini0879-matchessuppliedthreecapacitorfigure}
  \lean{PhyXMiniProblems.ProblemPhyXMini0879.MatchesSuppliedThreeCapacitorFigure}
  \uses{def:physics:phyx-mini-0879:phyxminiproblems-problemphyxmini0879-potentialdifferenceinvolts, def:physics:phyx-mini-0879:phyxminiproblems-problemphyxmini0879-capacitanceinmicrofarads, def:physics:phyx-mini-0879:phyxminiproblems-problemphyxmini0879-threeseriescapacitorsetup}
  Primary-raster evidence and calibration of its printed values to independent physical quantities. No charge value is printed in the image or included in these assumptions
\end{definition}
\begin{proof}
  This is the declared model definition of Matches Supplied Three Capacitor Figure.
\end{proof}
\begin{definition}[Has Physical Series Capacitor Parameters]
  \label{def:physics:phyx-mini-0879:phyxminiproblems-problemphyxmini0879-hasphysicalseriescapacitorparameters}
  \lean{PhyXMiniProblems.ProblemPhyXMini0879.HasPhysicalSeriesCapacitorParameters}
  \uses{def:physics:phyx-mini-0879:phyxminiproblems-problemphyxmini0879-potentialdifferenceinvolts, def:physics:phyx-mini-0879:phyxminiproblems-problemphyxmini0879-capacitanceinfarads, def:physics:phyx-mini-0879:phyxminiproblems-problemphyxmini0879-threeseriescapacitorsetup}
  Positivity and nondegeneracy conditions for the physical circuit
\end{definition}
\begin{proof}
  This is the declared model definition of Has Physical Series Capacitor Parameters.
\end{proof}
\begin{definition}[Satisfies Ideal Capacitor Laws]
  \label{def:physics:phyx-mini-0879:phyxminiproblems-problemphyxmini0879-satisfiesidealcapacitorlaws}
  \lean{PhyXMiniProblems.ProblemPhyXMini0879.SatisfiesIdealCapacitorLaws}
  \uses{def:physics:phyx-mini-0879:phyxminiproblems-problemphyxmini0879-potentialdifferenceinvolts, def:physics:phyx-mini-0879:phyxminiproblems-problemphyxmini0879-capacitanceinfarads, def:physics:phyx-mini-0879:phyxminiproblems-problemphyxmini0879-chargemagnitudeincoulombs, def:physics:phyx-mini-0879:phyxminiproblems-problemphyxmini0879-signedchargeincoulombs, def:physics:phyx-mini-0879:phyxminiproblems-problemphyxmini0879-threeseriescapacitorsetup}
  The ideal-capacitor constitutive law Q = C ΔV, together with the opposite signed charges on a capacitor's two plates. These laws apply uniformly to every capacitor and contain no circuit-specific numerical charge
\end{definition}
\begin{proof}
  This is the declared model definition of Satisfies Ideal Capacitor Laws.
\end{proof}
\begin{definition}[net Plate Charge At Node In Coulombs]
  \label{def:physics:phyx-mini-0879:phyxminiproblems-problemphyxmini0879-netplatechargeatnodeincoulombs}
  \lean{PhyXMiniProblems.ProblemPhyXMini0879.netPlateChargeAtNodeInCoulombs}
  \uses{def:physics:phyx-mini-0879:phyxminiproblems-problemphyxmini0879-signedchargeincoulombs, def:physics:phyx-mini-0879:phyxminiproblems-problemphyxmini0879-capacitorlabel, def:physics:phyx-mini-0879:phyxminiproblems-problemphyxmini0879-circuitnode, def:physics:phyx-mini-0879:phyxminiproblems-problemphyxmini0879-threeseriescapacitorsetup}
  Net signed charge residing on capacitor plates incident to one node
\end{definition}
\begin{proof}
  This is the declared model definition of net Plate Charge At Node In Coulombs.
\end{proof}
\begin{definition}[Is Floating Junction]
  \label{def:physics:phyx-mini-0879:phyxminiproblems-problemphyxmini0879-isfloatingjunction}
  \lean{PhyXMiniProblems.ProblemPhyXMini0879.IsFloatingJunction}
  \uses{def:physics:phyx-mini-0879:phyxminiproblems-problemphyxmini0879-circuitnode}
  The two internal nodes that have no direct connection to the source
\end{definition}
\begin{proof}
  This is the declared model definition of Is Floating Junction.
\end{proof}
\begin{definition}[Satisfies Series Circuit Conservation Laws]
  \label{def:physics:phyx-mini-0879:phyxminiproblems-problemphyxmini0879-satisfiesseriescircuitconservationlaws}
  \lean{PhyXMiniProblems.ProblemPhyXMini0879.SatisfiesSeriesCircuitConservationLaws}
  \uses{def:physics:phyx-mini-0879:phyxminiproblems-problemphyxmini0879-potentialdifferenceinvolts, def:physics:phyx-mini-0879:phyxminiproblems-problemphyxmini0879-capacitorlabel, def:physics:phyx-mini-0879:phyxminiproblems-problemphyxmini0879-threeseriescapacitorsetup, def:physics:phyx-mini-0879:phyxminiproblems-problemphyxmini0879-netplatechargeatnodeincoulombs, def:physics:phyx-mini-0879:phyxminiproblems-problemphyxmini0879-isfloatingjunction}
  Kirchhoff voltage balance around the closed loop and conservation of charge at every floating junction. Neutrality is stated at nodes rather than as an assumed equality of the three requested capacitor charges
\end{definition}
\begin{proof}
  This is the declared model definition of Satisfies Series Circuit Conservation Laws.
\end{proof}
\begin{lemma}[charge Magnitudes Are Equal]
  \label{lem:physics:phyx-mini-0879:phyxminiproblems-problemphyxmini0879-chargemagnitudesareequal}
  \lean{PhyXMiniProblems.ProblemPhyXMini0879.chargeMagnitudesAreEqual}
  \uses{def:physics:phyx-mini-0879:phyxminiproblems-problemphyxmini0879-chargemagnitudeincoulombs, def:physics:phyx-mini-0879:phyxminiproblems-problemphyxmini0879-threeseriescapacitorsetup, def:physics:phyx-mini-0879:phyxminiproblems-problemphyxmini0879-matchesthreecapacitorseriestopology, def:physics:phyx-mini-0879:phyxminiproblems-problemphyxmini0879-satisfiesidealcapacitorlaws, def:physics:phyx-mini-0879:phyxminiproblems-problemphyxmini0879-satisfiesseriescircuitconservationlaws}
  The two neutral floating junctions and opposite plate signs imply that all three capacitor charge magnitudes agree. This is derived from topology and general conservation rather than included in a premise
\end{lemma}
\begin{proof}
  The conclusion follows from the governing relations and figure readouts named in the statement of charge Magnitudes Are Equal.
\end{proof}
\begin{definition}[Answer Choice]
  \label{def:physics:phyx-mini-0879:phyxminiproblems-problemphyxmini0879-answerchoice}
  \lean{PhyXMiniProblems.ProblemPhyXMini0879.AnswerChoice}
  The four answer labels printed with the problem
\end{definition}
\begin{proof}
  This is the declared model definition of Answer Choice.
\end{proof}
\begin{definition}[displayed Charge In Microcoulombs]
  \label{def:physics:phyx-mini-0879:phyxminiproblems-problemphyxmini0879-answerchoice-displayedchargeinmicrocoulombs}
  \lean{PhyXMiniProblems.ProblemPhyXMini0879.AnswerChoice.displayedChargeInMicrocoulombs}
  \uses{def:physics:phyx-mini-0879:phyxminiproblems-problemphyxmini0879-answerchoice}
  The displayed common-charge magnitude for each choice, in microcoulombs
\end{definition}
\begin{proof}
  This is the declared model definition of displayed Charge In Microcoulombs.
\end{proof}
\begin{definition}[recorded Dataset Answer]
  \label{def:physics:phyx-mini-0879:phyxminiproblems-problemphyxmini0879-recordeddatasetanswer}
  \lean{PhyXMiniProblems.ProblemPhyXMini0879.recordedDatasetAnswer}
  \uses{def:physics:phyx-mini-0879:phyxminiproblems-problemphyxmini0879-answerchoice}
  The source dataset's recorded answer, retained as answer metadata
\end{definition}
\begin{proof}
  This is the declared model definition of recorded Dataset Answer.
\end{proof}
\begin{definition}[Answer Matches Common Charge]
  \label{def:physics:phyx-mini-0879:phyxminiproblems-problemphyxmini0879-answermatchescommoncharge}
  \lean{PhyXMiniProblems.ProblemPhyXMini0879.AnswerMatchesCommonCharge}
  \uses{def:physics:phyx-mini-0879:phyxminiproblems-problemphyxmini0879-chargemagnitudeinmicrocoulombs, def:physics:phyx-mini-0879:phyxminiproblems-problemphyxmini0879-threeseriescapacitorsetup, def:physics:phyx-mini-0879:phyxminiproblems-problemphyxmini0879-answerchoice}
  A displayed choice agrees with every physical capacitor charge
\end{definition}
\begin{proof}
  This is the declared model definition of Answer Matches Common Charge.
\end{proof}
\begin{definition}[Is Unique Matching Answer]
  \label{def:physics:phyx-mini-0879:phyxminiproblems-problemphyxmini0879-isuniquematchinganswer}
  \lean{PhyXMiniProblems.ProblemPhyXMini0879.IsUniqueMatchingAnswer}
  \uses{def:physics:phyx-mini-0879:phyxminiproblems-problemphyxmini0879-threeseriescapacitorsetup, def:physics:phyx-mini-0879:phyxminiproblems-problemphyxmini0879-answerchoice, def:physics:phyx-mini-0879:phyxminiproblems-problemphyxmini0879-answermatchescommoncharge}
  Exactly one displayed choice agrees with the common physical charge
\end{definition}
\begin{proof}
  This is the declared model definition of Is Unique Matching Answer.
\end{proof}
% --- Archon declaration topology end ---
% --- Archon physics formalization source end ---
